\scanpage{113}
\mspage{100}
% archive sheet number: 112

Ce couple satisfait aux conditions de Quillen
pour les couples de foncteurs adjoints entre
catégories de Quillen (Hom. Alg. Chap I,
\S~4, \unsure{th.|Th.} 3 p.~4.5), de sorte qu'ils
définissent par passage aux localisés un
couple de foncteurs adjoints
\diagram{\begin{tikzcd}[ampersand replacement=\&,column sep=large]
{(6.1.9)} \& \mathrm{Ho}_{\unsure{W|w}}(I) \arrow[r,shift left=1.25ex,"{u_{!}^{W}}"] \& \mathrm{Ho}_{W}(I') \arrow[l,shift left=1.25ex,"{u^{*}_{W}}"]
\end{tikzcd}}\unsure{,|.}
On a, plus précisément, les relations
\[
\left\{\begin{aligned}
 u_{!}(C(I))&\subset C(I')\\
 u_{!}(TC(I))&\subset TC(I')
\end{aligned}\right.\tag*{(6.1.10)}
\]
d'où
\[
 u_{!}(M(I)_{c})\subset M(I')_{c},\tag*{(6.1.11)}
\]
et de plus
\[
\begin{aligned}
 &\text{pour toute \unsure{flèches|flèche} }\inserted{f}\ \text{dans }M(I)_{c},\text{ telle}\\
 &\text{que }f\in W(I),\ \text{on a }\ u_{!}(f)\in W(I').
\end{aligned}\tag*{(6.1.12)}
\]
\inserted{i.e. $(u_{!})_{\unsure{c}}:M(I)_{c}\mathrel{-}M(I')_{c}$ est comp. avec les \unsure{localiseurs|localisations}}
D'autre part on a
\[
\left\{\begin{aligned}
 u^{*}(F(I'))&\subset F(I)\\
 u^{\unsure{*|4}}(TF(I'))&\subset TF(I)\\
 u^{*}(W(I'))&\subset W(I)
\end{aligned}\right.\tag*{(6.1.13)}
\]
donc
\[
\left\{\begin{aligned}
 &u^{*}(M(I')_{f})\subset M(I)_{f}\\
 &u^{*}_{f}:M(I')_{f}\mathrel{\unsure{-|\longrightarrow}}M(I)_{f}\ \text{est compatible}\\
 &\hspace{5em}\text{avec les localiseurs}.
\end{aligned}\right.\tag*{(6.1.14)}
\]

